%%%%%%%%%%%%%%%%%%%%%%%%%%%%%%%%%%%%%%%%%%%%%%%%%%%%%%%%%%%%%%%%%%%%%%%%%%%%%%%
%%
%%  LOD-macro.tex --- Macro File for the Lectures on Diffeology
%%
%%  Patrick Iglesias-Zemmour 2019-23
%%
%%%%%%%%%%%%%%%%%%%%%%%%%%%%%%%%%%%%%%%%%%%%%%%%%%%%%%%%%%%%%%%%%%%%%%%%%%%%%%%

%%%%%%%%%%%%%%%%%%%%%%%%%%%%%%%%%%%%%%%%%%%%%%%%%%%%%%%%%%%%%%%%%%%%%%%%%%%%%%%
%%
%% MARK: Text Shortcut
%%
%%%%%%%%%%%%%%%%%%%%%%%%%%%%%%%%%%%%%%%%%%%%%%%%%%%%%%%%%%%%%%%%%%%%%%%%%%%%%%%

%% Implement text inside a display math block to make space around the text
\newcommand{\SpacedText}[1]{\ \text{#1}\ }
\newcommand{\QuadText}[1]{\quad\text{#1}\quad}

%%%%%%%%%%%%%%%%%%%%%%%%%%%%%%%%%%%%%%%%%%%%%%%%%%%%%%%%%%%%%%%%%%%%%%%%%%%%%%%
%%
%% MARK: Hyphenation
%%
%%%%%%%%%%%%%%%%%%%%%%%%%%%%%%%%%%%%%%%%%%%%%%%%%%%%%%%%%%%%%%%%%%%%%%%%%%%%%%%

\hyphenation{pa-ra-me-t-ri-za-tion}
\hyphenation{pa-ra-me-t-ri-za-tions}
\hyphenation{dif-feo-lo-gi-cal}
\hyphenation{dif-feo-lo-gy}

%%%%%%%%%%%%%%%%%%%%%%%%%%%%%%%%%%%%%%%%%%%%%%%%%%%%%%%%%%%%%%%%%%%%%%%%%%%%%%%
%%
%% MARK: Latin shortcuts
%%
%%%%%%%%%%%%%%%%%%%%%%%%%%%%%%%%%%%%%%%%%%%%%%%%%%%%%%%%%%%%%%%%%%%%%%%%%%%%%%%

\newcommand{\cf}{{\em cf.\/}}
\newcommand{\ie}{{\em i.e.\/}}
\newcommand{\resp}{{\em resp.\/}}

%%%%%%%%%%%%%%%%%%%%%%%%%%%%%%%%%%%%%%%%%%%%%%%%%%%%%%%%%%%%%%%%%%%%%%%%%%%%%%%
%%
%% Define *** MARK: -
%%
%%%%%%%%%%%%%%%%%%%%%%%%%%%%%%%%%%%%%%%%%%%%%%%%%%%%%%%%%%%%%%%%%%%%%%%%%%%%%%%

%%%%%%%%%%%%%%%%%%%%%%%%%%%%%%%%%%%%%%%%%%%%%%%%%%%%%%%%%%%%%%%%%%%%%%%%%%%%%%%
%%
%% MARK: Numbers Sets
%%
%%%%%%%%%%%%%%%%%%%%%%%%%%%%%%%%%%%%%%%%%%%%%%%%%%%%%%%%%%%%%%%%%%%%%%%%%%%%%%%
 
 \providecommand{\NumberSetBold}[1]{\mathbf{#1}} %% Set the numbers set bold face

\newcommand{\RR}{\NumberSetBold{R}} %% Real numbers field
\newcommand{\CC}{\NumberSetBold{C}} %% Complex numbers field
\newcommand{\ZZ}{\NumberSetBold{Z}} %% Relative integers ring
\newcommand{\NN}{\NumberSetBold{N}} %% Natural integers
\newcommand{\QQ}{\NumberSetBold{Q}} %% Rational numbers field

%\newcommand{\fRR}{\boldsymbol{\mathsf R}} %% This command if we want to use a bold R sans-serif

%%%%%%%%%%%%%%%%%%%%%%%%%%%%%%%%%%%%%%%%%%%%%%%%%%%%%%%%%%%%%%%%%%%%%%%%%%%%%%%
%%
%% MARK: Various Arrows
%%
%%%%%%%%%%%%%%%%%%%%%%%%%%%%%%%%%%%%%%%%%%%%%%%%%%%%%%%%%%%%%%%%%%%%%%%%%%%%%%%

%% Attention ne pas toucher la longueur, affecte les diagrames / Changer le code pour utiliser tikz
\newcommand{\rfl}[1]{{\buildrel{\displaystyle #1}\over{\hbox to 9mm{\ \rightarrowfill \ }}}}

%%%%%%%%%%%%%%%%%%%%%%%%%%%%%%%%%%%%%%%%%%%%%%%%%%%%%%%%%%%%%%%%%%%%%%%%%%%%%%%
%%
%% Define *** MARK: -
%%
%%%%%%%%%%%%%%%%%%%%%%%%%%%%%%%%%%%%%%%%%%%%%%%%%%%%%%%%%%%%%%%%%%%%%%%%%%%%%%%

%%%%%%%%%%%%%%%%%%%%%%%%%%%%%%%%%%%%%%%%%%%%%%%%%%%%%%%%%%%%%%%%%%%%%%%%%%%%%%%
%%
%% MARK: Categories
%%
%%%%%%%%%%%%%%%%%%%%%%%%%%%%%%%%%%%%%%%%%%%%%%%%%%%%%%%%%%%%%%%%%%%%%%%%%%%%%%%

\newcommand{\Category}[1]{\{\text{#1}\}}  %% Display the name of the category between braces

\newcommand{\Obj}{{\operatorname{Obj}}}  %% The set of objects of a category
\newcommand{\Mor}{{\operatorname{Mor}}}  %% The set of Morphisms of a category
 
\newcommand{\Top}{\{\operatorname{Top}\}}         %% The sets of topologies of a set as in $\Top(\RR^n)$
\newcommand{\Difflg}{\operatorname{Diffeologies}} %% The set of all diffeologies on a set

%% Maps in categories
\newcommand{\Hom}{\mathrm{Hom}} %% Homomorphisms
\newcommand{\DHom}{\Hom^\infty} %% Smooth homomorphisms

\newcommand{\Isom}{\mathrm{Isom}} %% Isomorphisms
\newcommand{\DIsom}{\Isom^\infty} %% Smooth isomorphisms

%%%%%%%%%%%%%%%%%%%%%%%%%%%%%%%%%%%%%%%%%%%%%%%%%%%%%%%%%%%%%%%%%%%%%%%%%%%%%%%
%%
%% MARK: Intervals
%%
%%%%%%%%%%%%%%%%%%%%%%%%%%%%%%%%%%%%%%%%%%%%%%%%%%%%%%%%%%%%%%%%%%%%%%%%%%%%%%%

\newcommand{\OpenInterval}[1]{\mathopen] #1 \mathclose[}        %% An open interval both sides
\newcommand{\ClosedInterval}[1]{\mathopen[ #1 \mathclose]}      %% A closed interval both sides
\newcommand{\LeftClosedInterval}[1]{\mathopen[ #1 \mathclose[}  %% An interval closed left open right
\newcommand{\RightClosedInterval}[1]{\mathopen] #1 \mathclose]} %% An interval open left closed right

%%%%%%%%%%%%%%%%%%%%%%%%%%%%%%%%%%%%%%%%%%%%%%%%%%%%%%%%%%%%%%%%%%%%%%%%%%%%%%%
%%
%% MARK: Pseudo Groups
%%
%%%%%%%%%%%%%%%%%%%%%%%%%%%%%%%%%%%%%%%%%%%%%%%%%%%%%%%%%%%%%%%%%%%%%%%%%%%%%%%

\newcommand{\loc}{\mathrm{loc}}          %% Shortcut for "local". Used in $\Diff_\loc$ to denote th local diffeomorphisms etc.

\newcommand{\germ}{\operatorname{germ}}  %% The germ of a function

%%%%%%%%%%%%%%%%%%%%%%%%%%%%%%%%%%%%%%%%%%%%%%%%%%%%%%%%%%%%%%%%%%%%%%%%%%%%%%%
%%
%% MARK: Linear Forms
%%
%%%%%%%%%%%%%%%%%%%%%%%%%%%%%%%%%%%%%%%%%%%%%%%%%%%%%%%%%%%%%%%%%%%%%%%%%%%%%%%

\newcommand{\LinearForm}{\Lambda}                 %% Linear forms on Euclidean spaces, use as $\Form^p(E)
\newcommand{\SmoothForms}{{\pmb{\Lambda}}}        %% Smooth form space: on Euclidean domains

%%%%%%%%%%%%%%%%%%%%%%%%%%%%%%%%%%%%%%%%%%%%%%%%%%%%%%%%%%%%%%%%%%%%%%%%%%%%%%%
%%
%% MARK: Differential Calculus
%%
%%%%%%%%%%%%%%%%%%%%%%%%%%%%%%%%%%%%%%%%%%%%%%%%%%%%%%%%%%%%%%%%%%%%%%%%%%%%%%%

\newcommand{\Czero}{{\cC}^0}        %% Set of continuous maps
\newcommand{\Ck}[1]{{\cC}^{#1}}     %% Set of k-continuously differentiable maps
\newcommand{\Cinfty}{{\cC}^\infty}  %% Set of infinitely continuously differentiable maps

\newcommand{\Der}{\operatorname{D}}   %% Tangent linear map, use $\Der(F)_x(v)$ the tangent linear map of $F$ at a point $x$ applied to a vector $v$

\newcommand{\DForms}{\Omega} %% Differential forms, use as $\DForms^p(X)$

\newcommand{\Surf}{{\operatorname{Surf}}} %% Surface element - Volume element on a 2-dimensional manifold
\newcommand{\Vol}{\operatorname{Vol}}     %% Volume element

\newcommand{\CHO}{\mathsf{K}}                %% Chain-Homotopy operator
\newcommand{\Komega}{\CHO\omega}             %% SHortcut for the Action 1-form on the space of paths
\newcommand{\bfCHO}{\boldsymbol{\mathsf K}}  %% alternative to Chain-Homotopy operator

\newcommand{\dt}{\mathrm{dt}}  %% Differential symbol
\newcommand{\ds}{\mathrm{ds}}  %% Differential symbol
\newcommand{\du}{\mathrm{du}}  %% Differential symbol

%%%%%%%%%%%%%%%%%%%%%%%%%%%%%%%%%%%%%%%%%%%%%%%%%%%%%%%%%%%%%%%%%%%%%%%%%%%%%%%%
%%
%% MARK: Fiber Bundles
%%
%%%%%%%%%%%%%%%%%%%%%%%%%%%%%%%%%%%%%%%%%%%%%%%%%%%%%%%%%%%%%%%%%%%%%%%%%%%%%%%%

%% The set or space of section of any projection map
\newcommand{\Sections}{\operatorname{Sec}}

%%%%%%%%%%%%%%%%%%%%%%%%%%%%%%%%%%%%%%%%%%%%%%%%%%%%%%%%%%%%%%%%%%%%%%%%%%%%%%%
%%
%% MARK: Functions and Operators
%%
%%%%%%%%%%%%%%%%%%%%%%%%%%%%%%%%%%%%%%%%%%%%%%%%%%%%%%%%%%%%%%%%%%%%%%%%%%%%%%%

\newcommand{\Acos}{\mathop\mathrm{acos}\nolimits}  %% Function arc-cosinus
\newcommand{\Asin}{\mathop\mathrm{asin}\nolimits}  %% Function arc-sinus

\newcommand{\Const}{\mathrm{Const.}}               %% A constant function for a undefined constant

\newcommand{\Modulus}[1]{\vert #1 \vert}           %% The modulus of a complex number
\newcommand{\Norm}[1]{\Vert #1 \Vert}              %% The Euclidean norm of a vector

\newcommand{\Proj}{\operatorname{pr}}              %% Generic name for a projection

\newcommand{\Sign}{\operatorname{sgn}}             %% The function 'sign"

\renewcommand{\det}{\mathrm{det}}                   %% Determinant

\newcommand{\Value}{\operatorname{value}}           %% The value of something

\newcommand{\SquareRoot}{\operatorname{sqrt}}       %% Square root function

\newcommand{\Sup}{\operatorname{Sup}}               %% The supremum of a function

%%%%%%%%%%%%%%%%%%%%%%%%%%%%%%%%%%%%%%%%%%%%%%%%%%%%%%%%%%%%%%%%%%%%%%%%%%%%%%%
%%
%% MARK: Generating families
%%
%%%%%%%%%%%%%%%%%%%%%%%%%%%%%%%%%%%%%%%%%%%%%%%%%%%%%%%%%%%%%%%%%%%%%%%%%%%%%%%

\newcommand{\Gen}{\operatorname{Gen}}
\newcommand{\gen}[1]{\mathop\langle#1\mathop\rangle}

\newcommand{\vecspan}[1]{{\rm span}(#1)}

%%%%%%%%%%%%%%%%%%%%%%%%%%%%%%%%%%%%%%%%%%%%%%%%%%%%%%%%%%%%%%%%%%%%%%%%%%%%%%%
%%
%% MARK: Groups Names
%%
%%%%%%%%%%%%%%%%%%%%%%%%%%%%%%%%%%%%%%%%%%%%%%%%%%%%%%%%%%%%%%%%%%%%%%%%%%%%%%%

\newcommand{\GL}{\mathrm{GL}}     %% Linear group
\newcommand{\Sp}{{\rm Sp}}        %% Symplectic group
\newcommand{\SO}{\mathrm{SO}}     %% Special Orthogonal group
\newcommand{\SL}{\mathrm{SL}}     %% Special Linear group
\newcommand{\U}{\mathrm{U}}       %% Unitary group
\newcommand{\Perm}{\mathrm{Perm}} %% The group of permutations of a set
\newcommand{\SU}{\mathrm{SU}}

\let\O\relax
\newcommand{\O}{\mathrm{O}}       %% Orthogonal group

\newcommand{\Aff}{\operatorname{Aff}}   %% Affine group as in $\Aff(\RR)$
\newcommand{\Aut}{\operatorname{Aut}}   %% Automorphisms group
\newcommand{\Diff}{\operatorname{Diff}} %% Diffeomorphism group

%%%%%%%%%%%%%%%%%%%%%%%%%%%%%%%%%%%%%%%%%%%%%%%%%%%%%%%%%%%%%%%%%%%%%%%%%%%%%%%
%%
%% MARK: Groups Operators
%%
%%%%%%%%%%%%%%%%%%%%%%%%%%%%%%%%%%%%%%%%%%%%%%%%%%%%%%%%%%%%%%%%%%%%%%%%%%%%%%%

\newcommand{\coker}{\operatorname{coker}}   %% co-kernel
\newcommand{\Ad}{\mathrm{Ad}}               %% Adjoint action $\Ag(g)(k) = gkg^{-1}$

\DeclareMathOperator{\Lm}{L}    %% Left-multiplication in $\Lm(g)(k) = gk$
\DeclareMathOperator{\Rm}{R}    %% Right-multiplication in $\Rm(g)(k) = kg$

\newcommand{\OrbitMap}[1]{\hat{#1}}        %% Orbit map $\OrbitMap{x} : g \mapsto g(x)$

\newcommand{\Ab}[1]{#1^{\mathrm{Ab}}}      %% Abelianized, use as $\Ab{\pi_1}(X)$

\newcommand{\Stab}{\mathrm{St}}             %% Stabilizer

\DeclareMathOperator{\Trsl}{Tr} %% Translation operation, use as $\Trsl_a(x) = x + a$

%%%%%%%%%%%%%%%%%%%%%%%%%%%%%%%%%%%%%%%%%%%%%%%%%%%%%%%%%%%%%%%%%%%%%%%%%%%%%%%
%%
%% MARK: Lie Algebras
%%
%%%%%%%%%%%%%%%%%%%%%%%%%%%%%%%%%%%%%%%%%%%%%%%%%%%%%%%%%%%%%%%%%%%%%%%%%%%%%%%

%% Input the name of the group and it produces the name of the Lie algebra
%% as lower-case fracture font
\newcommand{\LieAlgebra}[1]{\mathfrak{\MakeLowercase{#1}}}

\newcommand{\DLie}{\mbox{\pounds}} %% Lie derivative

%%%%%%%%%%%%%%%%%%%%%%%%%%%%%%%%%%%%%%%%%%%%%%%%%%%%%%%%%%%%%%%%%%%%%%%%%%%%%%%
%%
%% MARK: Groupoids
%%
%%%%%%%%%%%%%%%%%%%%%%%%%%%%%%%%%%%%%%%%%%%%%%%%%%%%%%%%%%%%%%%%%%%%%%%%%%%%%%%

%% The source map
\newcommand{\Src}{\mathrm{src}}

%% The target map
\newcommand{\Trg}{\mathrm{trg}}

%% Transitivity component, also called orbit, in a groupoid. They are the largest subsets of objects connected by arrows
\newcommand{\TransComp}{\operatorname{Comp}}

%%%%%%%%%%%%%%%%%%%%%%%%%%%%%%%%%%%%%%%%%%%%%%%%%%%%%%%%%%%%%%%%%%%%%%%%%%%%%%%
%%
%% MARK: Homology - Cohomology
%%
%%%%%%%%%%%%%%%%%%%%%%%%%%%%%%%%%%%%%%%%%%%%%%%%%%%%%%%%%%%%%%%%%%%%%%%%%%%%%%%

%% The boundary operator in homology or anything realted
\DeclareMathOperator{\Boundary}{\partial}

\newcommand{\cub}[1]{\operatorname{I}^{#1}}      %% A cube, that is a product $I^p$ where $I$ is the interval $[0,1]$
\newcommand{\DCubes}[1]{\operatorname{Cub}_{#1}} %% The space of smooth cubes: $\DCubes{p}(X)$ is the space of smooth maps from $\RR^n$ to $X$.

\newcommand{\DChains}[1]{\operatorname{C}_{#1}}   %% Smooth chains in cubic homology: linear combination of smooth cubes.
\newcommand{\DCochains}[1]{\operatorname{C}^{#1}} %% Smooth cochains in cubic homology.

\newcommand{\QDChains}[1]{\operatorname{\textbf{\textsf C}}_{#1}}   %% Reduced smooth chain groups: quotient of the group of cubic chains by the subgroup of degenerate cubic chains.
\newcommand{\QDCochains}[1]{\operatorname{\textbf{\textsf C}}^{#1}} %% Smooth Cochains defined as morphisms from Reduced smooth chains (the first letter 'Q' is for 'Quotient').

\newcommand{\QDZ}{\operatorname{\textbf{\textsf Z}}} %% Cycle/cocycle notation depending of the place of the index: $\QDZ_p(X)$ cycles in the reduced cubic homology.
%% $\QDZ^p(X,\RR)$ cocycles in the reduced cubic cohomology with values in $\RR$.
\newcommand{\QDB}{\operatorname{\textbf{\textsf B}}} %% Boundary/coboundary notation depending of the place of the index: $\QDB_p(X)$ boundaries in the reduced cubic homology.
%% $\QDB^p(X,\RR)$ coboundaries in the reduced cubic cohomology with values in $\RR$.
\newcommand{\QDH}{\operatorname{\textbf{\textsf H}}} %% The reduced cubic homology or cohomology groups $\QDH = \QDZ/\QDB$.

\newcommand{\ZG}{Z} % generic Homology/Cohomology cycles.
\newcommand{\BG}{B} % generic Homology/Cohomology boundaries.
\newcommand{\HG}{H} % generic Homology/Cohomology groups $\HG= \ZG/\BG$.

\newcommand{\ZDR}{Z_\mathrm{dR}} %% De Rham Cocycles
\newcommand{\BDR}{B_\mathrm{dR}} %% De Rham Coboundaries
\newcommand{\HDR}{H_\mathrm{dR}} %% De Rham cohomology groups

\newcommand{\hdR}{h_{\mathrm{dR}}} %% De Rham homomorphism, takes the degree as exponent.

%%%%%%%%%%%%%%%%%%%%%%%%%%%%%%%%%%%%%%%%%%%%%%%%%%%%%%%%%%%%%%%%%%%%%%%%%%%%%%%
%%
%% MARK: Metric definitions and Covariant calculus
%%
%%%%%%%%%%%%%%%%%%%%%%%%%%%%%%%%%%%%%%%%%%%%%%%%%%%%%%%%%%%%%%%%%%%%%%%%%%%%%%%

\newcommand{\hdiv}{{\hat{\mathrm d}}\text{iv}} %% Covariant divergence with respect to some metric

\newcommand{\Length}{\mathrm{Length}}          %% The length od a path in the sense $\Length(\gamma) = \int_0^1 \sqrt{g(\gamma)_t(1)(1)} \dt$

%%%%%%%%%%%%%%%%%%%%%%%%%%%%%%%%%%%%%%%%%%%%%%%%%%%%%%%%%%%%%%%%%%%%%%%%%%%%%%%
%%
%% MARK: Homotopy
%%
%%%%%%%%%%%%%%%%%%%%%%%%%%%%%%%%%%%%%%%%%%%%%%%%%%%%%%%%%%%%%%%%%%%%%%%%%%%%%%%

\newcommand{\Poinc}[1]{\boldsymbol{#1}} %% The Poincaré groupod for any space in parameter.

%%%%%%%%%%%%%%%%%%%%%%%%%%%%%%%%%%%%%%%%%%%%%%%%%%%%%%%%%%%%%%%%%%%%%%%%%%%%%%%
%%
%% MARK: Linear maps and vector spaces
%%
%%%%%%%%%%%%%%%%%%%%%%%%%%%%%%%%%%%%%%%%%%%%%%%%%%%%%%%%%%%%%%%%%%%%%%%%%%%%%%%

\newcommand{\Lin}{\operatorname{L}} %% The set of linear maps between two vector spaces as in $\Lin(E,F)$
\newcommand{\DLin}{\Lin^\infty}     %% The set of smooth linear maps between two diffeological vector spaces

\newcommand{\Vector}[1]{\begin{pmatrix}#1\end{pmatrix}}              %% A vector as column with curve parentheis, usage : \Vector{a \\ b \\...}
                                                                    %% Can be used to form a matrix by $\Vector{a & b \\ c & d }$
\newcommand{\SqVect}[1]{\left[\begin{matrix}#1\end{matrix} \right]}  %% The same but between square brackets

\DeclareMathOperator{\Matrix}{Mat}   %% the set of matrices, use as $\Matrix_p$ the space of p x p matrices

\renewcommand{\ker}{\mathrm{ker}}    %% The kernel of a linear map $\ker(M)$

%% Some special operators
\newcommand{\J}{\operatorname{J}}    %% The $\pi/2$ rotation matrix

%%%%%%%%%%%%%%%%%%%%%%%%%%%%%%%%%%%%%%%%%%%%%%%%%%%%%%%%%%%%%%%%%%%%%%%%%%%%%%%
%%
%% MARK: Curves, Lines, Geodesics etc.
%%
%%%%%%%%%%%%%%%%%%%%%%%%%%%%%%%%%%%%%%%%%%%%%%%%%%%%%%%%%%%%%%%%%%%%%%%%%%%%%%%

\newcommand{\Geod}{\operatorname{Geod}}

\newcommand{\Traj}{\operatorname{Traj}} %% The space of trajectories
\newcommand{\traj}{\mathrm{traj}}       %% Subscript to identify the space of trajectories of curves

%%%%%%%%%%%%%%%%%%%%%%%%%%%%%%%%%%%%%%%%%%%%%%%%%%%%%%%%%%%%%%%%%%%%%%%%%%%%%%%
%%
%% MARK: Paths, Homotopy...
%%
%%%%%%%%%%%%%%%%%%%%%%%%%%%%%%%%%%%%%%%%%%%%%%%%%%%%%%%%%%%%%%%%%%%%%%%%%%%%%%%

\newcommand{\Paths}{\operatorname{Paths}}     %% Space of smooth paths in a diffeological space, use $\Paths(X).
\newcommand{\Loops}{\operatorname{Loops}}     %% The space of loops $\ell$, $\ell(0) = \ell(1)$

\newcommand{\stPaths}{\operatorname{stPaths}} %% Space of stationary paths in a diffeological space, use $st\Paths(X).
\newcommand{\stLoops}{\operatorname{stLoops}} %% Space of stationary loops in a diffeological space, use $\stLoops(X).

\newcommand{\0}{\hat 0}                       %% Origin of a path $\0(\gamma) = \gamma(0)$
\newcommand{\1}{\hat 1}                       %% End of a path $\1(\gamma) = \gamma(1)$

\def\Ends{\operatorname{ends}}                %% Ends of a path $\Ends(\gamma) = (\0(\gamma), \1(\gamma))$

\def\Rev{\operatorname{rev}}                  %% The reverse paths $\Rev(\gamma) = \bar\gamma = [t \mapsto \gamma(1-t)]$.

\newcommand{\Comp}{\operatorname{comp}}       %% Connected component map $\Comp(x)$ is the connected component containing $x$

%%%%%%%%%%%%%%%%%%%%%%%%%%%%%%%%%%%%%%%%%%%%%%%%%%%%%%%%%%%%%%%%%%%%%%%%%%%%%%%
%%
%% MARK: Diffeology
%%
%%%%%%%%%%%%%%%%%%%%%%%%%%%%%%%%%%%%%%%%%%%%%%%%%%%%%%%%%%%%%%%%%%%%%%%%%%%%%%%

\newcommand{\Domain}{\mathrm{dom}}       %% The domain of a parametrization
\newcommand{\Domains}{\mathrm{Domains}}  %% The set of Domains of an Euclidean space $\Domains(\RR^n)$

\newcommand{\Params}{\mathrm{Params}}  %% The set of parametrizations in a set
\newcommand{\Plots}{\mathrm{Plots}}    %% The set of plots in a diffeoogical space, the diffeology

%%%%%%%%%%%%%%%%%%%%%%%%%%%%%%%%%%%%%%%%%%%%%%%%%%%%%%%%%%%%%%%%%%%%%%%%%%%%%%%
%%
%% MARK: Various Maps
%%
%%%%%%%%%%%%%%%%%%%%%%%%%%%%%%%%%%%%%%%%%%%%%%%%%%%%%%%%%%%%%%%%%%%%%%%%%%%%%%%

\newcommand{\Def}{\mathrm{def}}          %% Set of definition of a map partially defined in a set

\newcommand{\Val}{\operatorname{Val}}     %% The space of values of a map/function

\newcommand{\Maps}{\mathrm{Maps}}         %% The set of all maps from a set to another $\Maps(A,B)$

\newcommand{\Id}{{\mathbf{1}}}            %% The identity map

\newcommand{\Supp}{\operatorname{Supp}}   %% The support of a function as in $\Supp(f) ± \{ x \mid f(x) \neq 0 \}$

\newcommand{\emptymap}{{\Id_\varnothing}} %% The empty map: a map being a subset of a product the empty map is the empty subset of the product

\newcommand{\Eval}{{\operatorname{ev}}}   %% Evaluation map : $\Eval(f,x) = f(x)$

\newcommand{\Rank}{\operatorname{rank}}   %% The rank of a linear map

\newcommand{\Sq}{{\operatorname{sq}}}     %% The square norm operator $\Sq(x) = \Norm{x}^2$

%%%%%%%%%%%%%%%%%%%%%%%%%%%%%%%%%%%%%%%%%%%%%%%%%%%%%%%%%%%%%%%%%%%%%%%%%%%%%%%
%%
%% MARK: Standard spaces
%%
%%%%%%%%%%%%%%%%%%%%%%%%%%%%%%%%%%%%%%%%%%%%%%%%%%%%%%%%%%%%%%%%%%%%%%%%%%%%%%%

\newcommand{\Sph}{\mathrm{S}}      %% Sphere, use $\Sph^n$ for n-sphere
\newcommand{\Torus}{\mathrm{T}}    %% Name of the torus $\RR/\ZZ$ with an index $\Torus_\alpha$ when describing an irratonal torus,
                                   %% or $\Torus_u$ when it is the space of geodesics of the torus with director vector $u \in \Sph^1$.
\newcommand{\HalfLine}{\Delta}     %% Half-line [0,∞[
\newcommand{\Corner}{\mathrm{K}}   %% Corner

%%%%%%%%%%%%%%%%%%%%%%%%%%%%%%%%%%%%%%%%%%%%%%%%%%%%%%%%%%%%%%%%%%%%%%%%%%%%%%%
%%
%% MARK: Symplectic Diffeology
%%
%%%%%%%%%%%%%%%%%%%%%%%%%%%%%%%%%%%%%%%%%%%%%%%%%%%%%%%%%%%%%%%%%%%%%%%%%%%%%%%

\newcommand{\Ham}{\operatorname{Ham}}   %% Hamiltonian diffeomorphisms
\newcommand{\grad}{\operatorname{grad}} %% Gradient vector

\newcommand{\Can}{\operatorname{Can}}   %% Canonical transformation — Also \Sympl

%%%%%%%%%%%%%%%%%%%%%%%%%%%%%%%%%%%%%%%%%%%%%%%%%%%%%%%%%%%%%%%%%%%%%%%%%%%%%%%
%%
%% MARK: Tangents space
%%
%%%%%%%%%%%%%%%%%%%%%%%%%%%%%%%%%%%%%%%%%%%%%%%%%%%%%%%%%%%%%%%%%%%%%%%%%%%%%%%

\newcommand{\Tgt}{T}             % Tangent space For $\T_x(X)$ the tangent space at $x$ to the space $X$
\newcommand{\TgtBundle}{T}       % Use as $\TgtBundle M$ or \TgtBundle{(X \times Y)} more letters  — tangent bundle
\newcommand{\UnitBundle}{U}      % Use as $\UnitBundle M$ (one letter) or \UnitBundle{(X \times Y)} more letters  — unit bundle

\newcommand{\VectorField}{\operatorname{Vect}}  %% The vector space of vector fields on a manifold

%%%%%%%%%%%%%%%%%%%%%%%%%%%%%%%%%%%%%%%%%%%%%%%%%%%%%%%%%%%%%%%%%%%%%%%%%%%%%%%
%%
%% MARK: Torus & Group of Periods
%%
%%%%%%%%%%%%%%%%%%%%%%%%%%%%%%%%%%%%%%%%%%%%%%%%%%%%%%%%%%%%%%%%%%%%%%%%%%%%%%%

\newcommand{\Periods}{\operatorname{P}} %% The group of periods, use $\Periods_\omega$

\newcommand{\cTorus}{\cT} %% Calligraphic letter T for disctinctive infinite torus.

%%%%%%%%%%%%%%%%%%%%%%%%%%%%%%%%%%%%%%%%%%%%%%%%%%%%%%%%%%%%%%%%%%%%%%%%%%%%%%%
%%
%% MARK: Classes and equivalence relations
%%
%%%%%%%%%%%%%%%%%%%%%%%%%%%%%%%%%%%%%%%%%%%%%%%%%%%%%%%%%%%%%%%%%%%%%%%%%%%%%%%

\newcommand{\Class}{\operatorname{class}}  %% Canonical projection from a set to its set of equivalence classes
\newcommand{\Cl}[1]{\left[ #1 \right]}     %% Short for the class of an element: $\Cl[x] = \Class(x)$

%%%%%%%%%%%%%%%%%%%%%%%%%%%%%%%%%%%%%%%%%%%%%%%%%%%%%%%%%%%%%%%%%%%%%%%%%%%%%%%
%%
%% Typography MARK: -
%%
%%%%%%%%%%%%%%%%%%%%%%%%%%%%%%%%%%%%%%%%%%%%%%%%%%%%%%%%%%%%%%%%%%%%%%%%%%%%%%%

%%%%%%%%%%%%%%%%%%%%%%%%%%%%%%%%%%%%%%%%%%%%%%%%%%%%%%%%%%%%%%%%%%%%%%%%%%%%%%%
%%
%% MARK: Bold characters
%%
%%%%%%%%%%%%%%%%%%%%%%%%%%%%%%%%%%%%%%%%%%%%%%%%%%%%%%%%%%%%%%%%%%%%%%%%%%%%%%%

%% All bold characters here except the number sets describe in a differfent place

\renewcommand{\AA}{{\mathbf A}}
\newcommand{\BB}{{\mathbf B}}
\newcommand{\DD}{{\mathbf D}}
\newcommand{\EE}{{\mathbf E}}
\newcommand{\FF}{{\mathbf F}}
\newcommand{\GG}{{\mathbf G}}
\newcommand{\HH}{{\mathbf H}}
\newcommand{\II}{{\mathbf I}}
\newcommand{\JJ}{{\mathbf J}}
\newcommand{\KK}{{\mathbf K}}
\newcommand{\LL}{{\mathbf L}}
\newcommand{\MM}{{\mathbf M}}
\newcommand{\OO}{{\mathbf O}}
\newcommand{\PP}{{\mathbf P}}
\renewcommand{\SS}{{\mathbf S}}
\newcommand{\TT}{{\mathbf T}}
\newcommand{\UU}{{\mathbf U}}
\newcommand{\VV}{{\mathbf V}}
\newcommand{\WW}{{\mathbf W}}
\newcommand{\XX}{{\mathbf X}}
\newcommand{\YY}{{\mathbf Y}}

\newcommand{\bmP}{\boldsymbol P}

\newcommand{\bmG}{{\mbox{\boldmath $G$}}}

\newcommand{\bPhi}{{\mbox{\boldmath $\Phi$}}}

\newcommand{\bS}{\fontshape{n}\textsf{\textbf S}}

\renewcommand{\aa}{{\bf a}}
\newcommand{\bb}{{\bf b}}
\newcommand{\ff}{{\bf f}}
\newcommand{\jj}{{\bf j}}
\newcommand{\rr}{{\bf r}}
\newcommand{\pp}{{\bf p}}
\renewcommand{\ss}{{\bf s}}
\newcommand{\uu}{{\bf u}}
\newcommand{\vv}{{\bf v}}
\newcommand{\xx}{{\bf x}}
\newcommand{\yy}{{\bf y}}
\newcommand{\ee}{{\bf e}}

\newcommand{\balpha}{\boldsymbol \alpha}
\newcommand{\bbeta}{\boldsymbol \beta}
\newcommand{\bomega}{\boldsymbol \omega}
\newcommand{\bOmega}{\boldsymbol \Omega}

\newcommand{\bgamma}{\boldsymbol \gamma}
\newcommand{\bvarphi}{\boldsymbol \varphi}
\newcommand{\blambda}{\boldsymbol \lambda}
\newcommand{\bsigma}{\boldsymbol \sigma}
\newcommand{\bpi}{\boldsymbol \pi}
\newcommand{\btau}{\boldsymbol \tau}

\newcommand{\bma}{{\boldsymbol a}}
\newcommand{\bmb}{{\boldsymbol b}}
\newcommand{\bmc}{{\boldsymbol c}}
\newcommand{\bme}{{\boldsymbol e}}
\newcommand{\bmf}{{\boldsymbol f}}
\newcommand{\bmg}{{\boldsymbol g}}
\newcommand{\bmo}{{\boldsymbol o}}
\newcommand{\bmp}{{\boldsymbol p}}
\newcommand{\bmr}{{\boldsymbol r}}
\newcommand{\bms}{{\boldsymbol s}}
\newcommand{\bmt}{{\boldsymbol t}}
\newcommand{\bmv}{{\boldsymbol v}}
\newcommand{\bmx}{{\boldsymbol x}}
\newcommand{\bmy}{{\boldsymbol y}}
\newcommand{\bmz}{{\boldsymbol z}}

\newcommand{\bmZero}{\boldsymbol 0}

\newcommand{\sbmx}{{\mbox{\boldmath $\scriptstyle x$}}}

\newcommand{\sbgamma}{{\mbox{\boldmath $\scriptstyle \gamma$}}}

%%%%%%%%%%%%%%%%%%%%%%%%%%%%%%%%%%%%%%%%%%%%%%%%%%%%%%%%%%%%%%%%%%%%%%%%%%%%%%%
%%
%% MARK: Math calligraphy
%%
%%%%%%%%%%%%%%%%%%%%%%%%%%%%%%%%%%%%%%%%%%%%%%%%%%%%%%%%%%%%%%%%%%%%%%%%%%%%%%%

\newcommand{\cA}{\mathcal{A}}
\newcommand{\cB}{\mathcal{B}}
\newcommand{\cC}{\mathcal{C}}
\newcommand{\cD}{\mathcal{D}}
\newcommand{\cE}{\mathcal{E}}
\newcommand{\cF}{\mathcal{F}}
\newcommand{\cG}{\mathcal{G}}
\newcommand{\cH}{\mathcal{H}}
\newcommand{\cI}{\mathcal{I}}
\newcommand{\cJ}{\mathcal{J}}
\newcommand{\cK}{\mathcal{K}}
\newcommand{\cL}{\mathcal{L}}
\newcommand{\cM}{\mathcal{M}}
\newcommand{\cN}{\mathcal{N}}
\newcommand{\cO}{\mathcal{O}}
\newcommand{\cP}{\mathcal{P}}
\newcommand{\cQ}{\mathcal{Q}}
\newcommand{\cR}{\mathcal{R}}
\newcommand{\cS}{\mathcal{S}}
\newcommand{\cT}{\mathcal{T}}
\newcommand{\cU}{\mathcal{U}}
\newcommand{\cV}{\mathcal{V}}
\newcommand{\cW}{\mathcal{W}}
\newcommand{\cX}{\mathcal{X}}
\newcommand{\cY}{\mathcal{Y}}
\newcommand{\cZ}{\mathcal{Z}}

%%%%%%%%%%%%%%%%%%%%%%%%%%%%%%%%%%%%%%%%%%%%%%%%%%%%%%%%%%%%%%%%%%%%%%%%%%%%%%%
%%
%% MARK: Math frak
%%
%%%%%%%%%%%%%%%%%%%%%%%%%%%%%%%%%%%%%%%%%%%%%%%%%%%%%%%%%%%%%%%%%%%%%%%%%%%%%%%

\newcommand{\fA}{{\mathfrak A}}
\newcommand{\fB}{{\mathfrak B}}
\newcommand{\fC}{{\mathfrak C}}
\newcommand{\fD}{{\mathfrak D}}
\newcommand{\fF}{{\mathfrak F}}
\newcommand{\fG}{{\mathfrak G}}
\newcommand{\fK}{{\mathfrak K}}
\newcommand{\fI}{{\mathfrak I}}
\newcommand{\fL}{{\mathfrak L}}
\newcommand{\fP}{{\mathfrak P}}
\newcommand{\fR}{{\mathfrak R}}
\newcommand{\fS}{{\mathfrak S}}
\newcommand{\fU}{{\mathfrak U}}
\newcommand{\fX}{{\mathfrak X}}
\newcommand{\fY}{{\mathfrak Y}}
\newcommand{\fZ}{{\mathfrak Z}}

\newcommand{\fu}{{\mathfrak u}}

